\documentclass[11pt]{article}
\usepackage{amsmath,amssymb,amsthm,hyperref}
\newtheorem{theorem}{Theorem}
\newtheorem{lemma}[theorem]{Lemma}
\newtheorem{corollary}[theorem]{Corollary}
\newtheorem{remark}[theorem]{Remark}
\newcommand{\E}{\mathbb E}
\begin{document}
\title{A monotone comparison-tensor obstruction for a small exceptional die}
\author{Research notes, 30 September 2026}
\maketitle

\paragraph{Status.} The structural theorem below has a complete elementary
proof. Together with the independently checked arithmetic result in the
companion note, it gives a universal $n-1$ per-die lower bound. This bound
is attained for $n=4$ by the separately certified construction.
No manuscript source was edited.

Let $D_0$ have $c$ faces in an $(m+1)$-die permutation-fair collection.
For its faces $q=1,\ldots,c$, in increasing order, put
$p_j(q)=\Pr(D_j<D_0(q))$, $1\leq j\leq m$.  These are nondecreasing
rational-valued functions of $q$.  Write $\E$ for uniform averaging over $q$.
Permutation fairness on every subset gives
\begin{equation}\label{eq:moments}
 \E\prod_{j\in S}p_j=\frac1{|S|+1}
 \qquad(S\subseteq[m]).
\end{equation}
Only these moments and the common monotone order are needed below.

\begin{theorem}[Reduction of a short monotone mixture]\label{thm:reduction}
Suppose $p_1,\ldots,p_m:\{1,\ldots,c\}\to[0,1]$ are nondecreasing
and satisfy \eqref{eq:moments}. If $c<m$, there is an index $j_0$ such that
\[
 \frac1c\sum_{q=1}^c p_{j_0}(q)^a=\frac1{a+1}
 \qquad(0\leq a\leq m).
\]
In particular, the $c$ numbers $p_{j_0}(q)$ form an equal-weight interval
quadrature of degree $m$. They are rational if all $p_j(q)$ are rational.
\end{theorem}

\begin{proof}
For real-valued functions on the $c$-point probability space, use the inner
product $\langle f,g\rangle=\E(fg)$.  Let $v_j=p_j-1/2$ and
$\rho=1/12$.  The first two orders of \eqref{eq:moments} give
$\E v_j=0$ and $\langle v_i,v_j\rangle=\rho$ for $i\ne j$.
Thus the Gram matrix has the form
\[
 G=D+\rho\boldsymbol 1\boldsymbol 1^T,\qquad
 D=\operatorname{diag}(d_1,\ldots,d_m),\quad
 d_j=\lVert v_j\rVert^2-\rho.
\]
All $v_j$ lie in a $(c-1)$-dimensional space, so $\operatorname{rank}G\leq c-1$.
Consequently $\operatorname{rank}D\leq c$. Since $m>c$, some $d_{j_0}=0$.
Set $p=p_{j_0}$ and $v=p-1/2$.
For every $j\ne j_0$, the function $e_j=p_j-p$ satisfies
\[
 \langle e_j,1\rangle=\langle e_j,v\rangle=0,
 \qquad \lVert e_j\rVert^2=d_j\geq0.
\]
Moreover $\langle e_i,e_j\rangle=0$ for distinct $i,j\ne j_0$.
Thus the nonzero $e_j$ are mutually orthogonal and are orthogonal to
both $1$ and the nonconstant $v$. There are at most $c-2$ of them.
Let $k$ be their number and put $r=m-k$. Exactly $r$ of the $p_j$ equal
$p$, and $r\geq m-c+2$.

Let $t$ be the number of distinct values taken by $p$.
Products of distinct members of the $r$-element equal group show that
\[
 \E p^a=\frac1{a+1}\quad(0\leq a\leq r).
\]
For an exceptional index $j$ and $0\leq a<r$, choosing $a$ distinct
members of that group gives
\[
 \E(p^a p_j)=\frac1{a+2}=\E p^{a+1},
 \qquad \E(p^a e_j)=0.
\]
If $t\geq r$, the $r$ functions $1,p,\ldots,p^{r-1}$ are linearly independent.
The $k$ nonzero orthogonal functions $e_j$ lie in their common orthogonal
complement, so $k\leq c-r$. This contradicts $k+r=m>c$.
Hence $t<r$.

Polynomial interpolation on the $t$ values of $p$ now gives
\begin{equation}\label{eq:condzero}
 \E(e_j\mid p)=0
\end{equation}
for every exceptional $j$. Since $p$ is nondecreasing, each level set of
$p$ is a consecutive block of the ordered $q$'s. On such a block $B$,
every $e_j=p_j-p$ is nondecreasing. Chebyshev's sum inequality and
\eqref{eq:condzero} therefore give
\[
 \E(e_i e_j\mid B)\geq0.
\]
Its unconditional expectation is zero by the already proved orthogonality,
so it is zero on every block. For two nondecreasing real sequences on a
finite interval, equality in this covariance inequality implies that at
least one sequence is constant: indeed,
\[
 \operatorname{Cov}_B(f,g)
 =\frac1{2|B|^2}\sum_{a,b\in B}(f(a)-f(b))(g(a)-g(b)),
\]
and if both sequences are nonconstant, the summand from the first and last
positions is strictly positive. By \eqref{eq:condzero}, a constant $e_j$
on $B$ is zero. Thus, on each level block of $p$, at most one exceptional
function $e_j$ is nonzero.

For any subset $S\subseteq[m]$, the product $\prod_{j\in S}p_j$ consequently
has, on a level block $p=x$, either the constant value $x^{|S|}$ or the
form $x^{|S|}+x^{|S|-1}e_j$. In either case its conditional expectation is
$x^{|S|}$. Hence
\[
 \E p^{|S|}=\E\prod_{j\in S}p_j=\frac1{|S|+1}.
\]
Taking sets of every size from $0$ to $m$ proves the theorem.
\end{proof}

\begin{remark}[Positive atom weights and repeated joint rows]
The structural proof also holds for arbitrary strictly positive weights
on the $c$ ordered atoms, with the conclusion $\E p^a=1/(a+1)$.
All dimension and interpolation arguments are unchanged. On a level
block, the only modification is the weighted covariance identity
\[
 \operatorname{Cov}(f,g)
 =\frac12\sum_{a,b}w_aw_b(f(a)-f(b))(g(a)-g(b)),
\]
where the positive conditional weights sum to one. Its equality case
for nondecreasing functions is the same as above.
Consequently, if a uniform comparison tensor has fewer than $m$
distinct joint rows, merging equal consecutive rows and then undoing
the merge gives a scalar quadrature on its original number of equal
atoms, including their repetitions.
\end{remark}

\begin{corollary}[Universal per-die bound]
Every die in every permutation-fair collection of $n\geq2$ dice has at
least $n-1$ faces.
\end{corollary}
\begin{proof}
Put $m=n-1$ and suppose that a die has $c<m$ faces. By
Theorem~\ref{thm:reduction}, there is a rational equal-weight interval
quadrature with $c$ nodes and degree $m$. The arithmetic theorem proved
in the companion note \texttt{rational\_designs/two\_adic\_lower\_bound.tex}
says that such a rule has at least $2^{\lfloor m/2\rfloor}$ nodes.
For $m\geq6$, this is at least $m$, a contradiction.
For $2\leq m\leq5$, only $c\leq4$ need be considered. A one-node rule
cannot have degree two. For $c=2,3,4$, map the nodes to $[-1,1]$.
Since $m>c$, their first $c$ power sums would equal those of the uniform
measure, multiplied by $c$. Newton identities force their node polynomial,
up to scalar, to be respectively
\[
 3x^2-1,\qquad x(2x^2-1),\qquad45x^4-30x^2+1.
\]
None splits into rational linear factors: the first two have visibly
irrational roots, and the squared roots of the last are
$(5\pm2\sqrt5)/15$. Thus no such rational quadrature exists. For $m=1$,
the desired bound is just the requirement that each die have a face.
\end{proof}

\paragraph{Independent classical alternative.}
Without the new arithmetic theorem, the structural result already gives
$c\geq n-1$ for all sufficiently large $n$: uniform-interval equal-weight
quadrature of degree $m$ requires $\Omega(m^2)$ nodes, contradicting $c<m$.
The classical Bernstein theorem and the finite calculations below can also
remove the remaining small cases.

\paragraph{Sources for the classical input.}
The lower quadratic order is stated in Kuijlaars,
\emph{The minimal number of nodes in Chebyshev type quadrature formulas}
(1993), DOI \href{https://doi.org/10.1016/0019-3577(93)90007-L}
{10.1016/0019-3577(93)90007-L}; it also follows from Gilboa--Peled,
\emph{Chebyshev-type quadratures for doubling weights},
\href{https://arxiv.org/abs/1507.01505}{arXiv:1507.01505}.
A stronger classical theorem of Bernstein says that an equal-weight
$c$-node uniform-interval quadrature of degree at least $c$ can have all
nodes real only for $c\in\{1,2,3,4,5,6,7,9\}$.
A primary paper explicitly recording this theorem is Gautschi--Monegato,
\emph{On Optimal Chebyshev-Type Quadratures}, Numerische Mathematik
28 (1977), 59--67,
\href{https://www.cs.purdue.edu/homes/wxg/selected_works/section_06/058.pdf}
{author-hosted PDF}. Its input is not needed if the parallel rational
quadrature lower bound is used.

\paragraph{Small rational cases.}
After mapping nodes to $[-1,1]$, a $c$-node degree-$c$ quadrature has power
sums $s_a=0$ for odd $a$ and $s_a=c/(a+1)$ for even $a\leq c$.
Newton identities uniquely determine the monic polynomial of its nodes.
For the exceptional $c$ permitted by Bernstein, the nontrivial factors are
\[
\begin{array}{c|l}
 c&\text{monic polynomial, up to a nonzero scalar}\\ \hline
 2&3x^2-1\\
 3&x(2x^2-1)\\
 4&45x^4-30x^2+1\\
 5&x(72x^4-60x^2+7)\\
 6&105x^6-105x^4+21x^2-1\\
 7&x(6480x^6-7560x^4+2142x^2-149)\\
 9&x(22400x^8-33600x^6+15120x^4-2280x^2+53).
\end{array}
\]
Each displayed nonconstant factor other than $x$ is irreducible over
$\mathbb Q$ (verified exactly by the accompanying SymPy script), so none
has all rational roots. For a short stand-alone proof, finite rational-root
tests suffice, and their complete integer certificates can be recorded.
Together with Bernstein, this rules out rational equal-weight degree-$m$
quadrature with $2\leq c\leq m$, implying $c\geq n-1$ for all $n\geq3$.
The case $c=1$, $m\geq2$ fails already by the first two moments.

\paragraph{Interpretation and remaining obstacle.}
The theorem does not prove an exponential per-die bound: it applies only
when $c<n-1$. The arithmetic rational-design obstruction, however strong,
therefore cannot simply be applied to every die. Extending the reduction
past this low-rank threshold would require new information. The existing
three-die example with a two-face die also shows that an unrestricted
rational-quadrature reduction is false: its two comparison functions give
a rational product mixture, but there is no rational equal-weight two-node
degree-two interval rule.
\end{document}
